%% toc: P15 170 170.10
% Proofs for simultaneous concentration bands, confidence sequences and their comparisons.
\section{Simultaneous concentration bands and confidence sequences: proofs}
\label{sec:annex-p15-w02b-p15-063}

\begin{annexitem}{annex:p15-w02b-p15-063-exponential-optimality}{thm:p15-w02b-p15-063-exponential-optimality}{Minimizing the exponential bound under the extremal law}
\begin{proof}
Write \(\mu_i=EX_i\), so \(\sum_i\mu_i=np\), and for real \(h\) put
\[
 \kappa(h)=-h(p+t)+\ln(1-p+pe^h).
\]

\emph{Part~1.} For \(h\leq0\), Jensen's inequality for the convex
exponential gives \(Ee^{hS_n}\geq e^{hES_n}=1\), hence
\(\Psi(h)\geq e^{-hnt}\geq1\). For \(h>0\), independence and the chord
bound of \cref{thm:p15-conc-exponential} on \([0,1]\) give
\[
 Ee^{hS_n}=e^{-hnp}\prod_iEe^{hX_i}
     \leq e^{-hnp}\prod_i(1-\mu_i+\mu_ie^h)
     \leq e^{-hnp}(1-p+pe^h)^n,
\]
the last step by the arithmetic-geometric mean inequality, as in the proof
of \cref{thm:p15-conc-common-range}. Thus \(\Psi(h)\leq e^{n\kappa(h)}\)
for \(h>0\). At \(h=\lambda_*\), computed below, \(e^{n\kappa(\lambda_*)}=
e^{-nD(p+t\Vert p)}\). The divergence is at least \(2t^2>0\) by the third
bound of \cref{thm:p15-conc-common-range} (its proof shows
\(D(p+t\Vert p)\geq g(p)t^2\geq2t^2\)), so this value is smaller than one.
Hence the infimum over \(h>0\) is smaller than one, which is at most every
value with \(h\leq0\); the two infima coincide and are at most
\(e^{-nD(p+t\Vert p)}\).

\emph{Part~2.} For Bernoulli variables of mean \(p\),
\(Ee^{hX_i}=1-p+pe^h\) exactly, so \(\Psi(h)=e^{n\kappa(h)}\) for every
real \(h\). Differentiating,
\[
 \kappa'(h)=-(p+t)+u_h,\quad
 \kappa''(h)=u_h(1-u_h)>0,\qquad u_h=\frac{pe^h}{1-p+pe^h}\in(0,1).
\]
So \(\kappa\) is strictly convex and its only critical point solves
\(u_h=p+t\), that is \(pe^h(1-p-t)=(p+t)(1-p)\), giving
\(h=\lambda_*\). Since \((p+t)(1-p)-p(1-p-t)=t>0\), \(\lambda_*>0\).
At \(\lambda_*\),
\[
 1-p+pe^{\lambda_*}=(1-p)+\frac{(p+t)(1-p)}{1-p-t}=\frac{1-p}{1-p-t},
\]
and therefore
\[
\begin{split}
 \kappa(\lambda_*)
 &=-(p+t)\ln\frac{p+t}{p}-(p+t)\ln\frac{1-p}{1-p-t}+\ln\frac{1-p}{1-p-t}\\
 &=-(p+t)\ln\frac{p+t}{p}-(1-p-t)\ln\frac{1-p-t}{1-p}
  =-D(p+t\Vert p).
\end{split}
\]
A strictly convex function attains its minimum only at its critical point,
so \(\inf_{h>0}\Psi(h)=e^{-nD(p+t\Vert p)}\), attained exactly at
\(\lambda_*\).

\emph{Part~3.} Independent Bernoulli variables of common mean \(p\) take
values in \([0,1]\) and have average mean \(p\), so they are admissible.
For them the infimum equals \(e^{-nD(p+t\Vert p)}\) by part~2, and the
assumed inequality at this law gives \(B(n,p,t)\geq e^{-nD(p+t\Vert p)}\).
By part~1, \(A_1\) itself satisfies the assumed inequality for every
admissible law; it is therefore the smallest such number.

\emph{Part~4.} Let \(q=p+t\in(p,1)\). For any law of \(X_1\) on \([0,1]\)
with mean \(p\), Markov's inequality for the nonnegative variable \(X_1\)
gives \(P(X_1\geq q)\leq p/q\). The two-point law with
\(P(X_1=q)=p/q\) and \(P(X_1=0)=1-p/q\) lies in \([0,1]\), has mean \(p\)
and attains this value, so the supremum is \(p/q\). Finally
\[
 \frac{e^{-D(q\Vert p)}}{p/q}
  =\left(\frac pq\right)^{q-1}\left(\frac{1-p}{1-q}\right)^{1-q}
  =\left(\frac qp\cdot\frac{1-p}{1-q}\right)^{1-q}>1,
\]
since both factors in the base exceed one and \(1-q>0\).
\end{proof}
\end{annexitem}

\begin{annexitem}{annex:p15-w02b-p15-063-chebyshev-comparison}{thm:p15-w02b-p15-063-chebyshev-comparison}{The variance ceiling and the crossing statements}
\begin{proof}
\emph{Part~1.} For \(0\leq X_i\leq1\), \(X_i^2\leq X_i\), so
\(\operatorname{Var}X_i=EX_i^2-\mu_i^2\leq\mu_i-\mu_i^2\), with equality
exactly when \(E(X_i-X_i^2)=0\). The variable \(X_i-X_i^2\) is nonnegative,
so equality means \(X_i^2=X_i\), that is \(X_i\in\{0,1\}\), almost surely.
Next,
\[
 \frac1n\sum_i\mu_i(1-\mu_i)=p-\frac1n\sum_i\mu_i^2
   =p-p^2-\frac1n\sum_i(\mu_i-p)^2\leq p(1-p),
\]
with equality exactly when every \(\mu_i=p\). Finally
\(p(1-p)=1/4-(p-1/2)^2\leq1/4\).

\emph{Part~2.} Since \(\sigma^2>0\), \(A_3<K_6\) is equivalent to
\(\sigma^2+x<\sigma^2e^{2x}\), that is to \(\varphi(x)>0\), where
\[
 \varphi(x)=\sigma^2e^{2x}-\sigma^2-x .
\]
Equalities and reverse inequalities correspond in the same way. Now
\(\varphi(0)=0\), \(\varphi'(x)=2\sigma^2e^{2x}-1\) and
\(\varphi''(x)=4\sigma^2e^{2x}>0\). By part~1, \(\varphi'(0)=2\sigma^2-1
\leq-1/2<0\). The derivative is strictly increasing and vanishes at
\(x_m=\tfrac12\ln(1/(2\sigma^2))>0\), so \(\varphi\) strictly decreases on
\([0,x_m]\) and strictly increases on \([x_m,\infty)\), with
\(\varphi(x)\geq\sigma^2(e^{2x}-1)-x\to\infty\). Hence \(\varphi<0\) on
\((0,x_m]\), and on \([x_m,\infty)\) there is exactly one zero \(x_*\),
with \(\varphi<0\) on \((0,x_*)\) and \(\varphi>0\) on \((x_*,\infty)\).

For the bound on \(x_*\), put \(x_1=1/(2\sigma^2)-1\geq1\). For \(x>0\)
the exponential series gives \(e^{2x}>1+2x+2x^2\), so
\[
 \varphi(x_1)>\sigma^2(2x_1+2x_1^2)-x_1
   =x_1\bigl(2\sigma^2(1+x_1)-1\bigr)=0 .
\]
Since \(x_1>0\) and \(\varphi(x_1)>0\), \(x_1>x_*\).
The chain \(A_1\leq A_2\leq A_3\) is \cref{thm:p15-conc-common-range},
and \(K_6\leq K_5\) because \(\sigma^2+nt^2\geq nt^2\).

\emph{Part~3.} By the expansion of \(G(t,p)=D(p+t\Vert p)/t^2\) in the
proof of \cref{thm:p15-conc-common-comparison}, valid for every
\(0<p<1\), \(G(t,p)\to1/[2p(1-p)]\) as \(t\downarrow0\). Since
\((1-e^{-y})/y\to1\) as \(y\downarrow0\),
\[
 \frac{1-A_1}{t^2}\longrightarrow\frac n{2p(1-p)},\qquad
 \frac{1-K_6}{t^2}=\frac n{\sigma^2+nt^2}\longrightarrow\frac n{\sigma^2}.
\]
By part~1, \(n/\sigma^2\geq n/[p(1-p)]>n/[2p(1-p)]\). The difference
\((A_1-K_6)/t^2=[(1-K_6)-(1-A_1)]/t^2\) therefore has a positive limit,
so \(A_1>K_6\) for all sufficiently small \(t>0\).

\emph{Part~4.} \(A_3=e^{-2nt^2}\) is an exponential in \(n\), and the
chain of part~2 bounds \(A_1,A_2\) by it. The function
\(v\mapsto v/(v+nt^2)\) is increasing in \(v>0\), so \(K_6\geq s/(s+nt^2)\).
Then
\[
 0<\frac{A_3}{K_6}\leq\frac{(s+nt^2)e^{-2nt^2}}{s}\longrightarrow0,
\]
because \(ne^{-2nt^2}\to0\) for fixed \(t>0\).

\emph{Part~5.} With \(a_i=0\), \(b_i=1\), \(V_n=n\), and the event
\(|\overline X-p|\geq t\) is \(|S_n|\geq nt\), so
\cref{thm:p15-conc-unequal-ranges} gives
\(P(|\overline X-p|\geq t)\leq2e^{-2n^2t^2/n}=2A_3\). Since \(x>0\),
\(2A_3<K_5\) means \(2e^{-2x}<\sigma^2/x\), that is \(\chi(x)<\sigma^2\)
with \(\chi(x)=2xe^{-2x}\); likewise \(2A_3\geq K_5\) means
\(\chi(x)\geq\sigma^2\). Now
\(\chi'(x)=2e^{-2x}(1-2x)\), so \(\chi\) is continuous, strictly
increasing on \((0,\tfrac12]\) and strictly decreasing on
\([\tfrac12,\infty)\), with \(\chi(x)\to0\) as \(x\downarrow0\) and as
\(x\to\infty\). At \(x=\tfrac12\ln2\), \(\chi=\tfrac12\ln2>\tfrac14\geq
\sigma^2\), using part~1. By the intermediate value theorem and strict
monotonicity there is exactly one \(x_-\in(0,\tfrac12\ln2)\) and exactly
one \(x_+\in(\tfrac12,\infty)\) with \(\chi=\sigma^2\), and
\(\chi\geq\sigma^2\) exactly on \([x_-,x_+]\).

For the upper limit put \(L=\ln(2/\sigma^2)\geq\ln8>\tfrac12\). Then
\[
 \chi(L)=2L\frac{\sigma^4}4=\sigma^2\cdot\frac{\sigma^2\ln(2/\sigma^2)}2 .
\]
The function \(s\mapsto s\ln(2/s)\) has derivative \(\ln(2/s)-1>0\) for
\(0<s<2/e\), and \(\tfrac14<2/e\), so for \(0<\sigma^2\leq\tfrac14\) the
last factor is at most \(\tfrac18\ln8<1\). Hence \(\chi(L)<\sigma^2\);
since \(L>\tfrac12\) and \(\chi\geq\sigma^2\) on \([\tfrac12,x_+]\),
\(L>x_+\).

Finally, for \(x<x_-<\tfrac12\ln2\), \(2A_3=2e^{-2x}>1\), and
\(K_5>2A_3\) because \(x\) lies outside \([x_-,x_+]\).
\end{proof}
\end{annexitem}

\begin{annexitem}{annex:p15-w02b-p15-063-normal-benchmark}{thm:p15-w02b-p15-063-normal-benchmark}{Gaussian optimization, variance ceilings and the Mills ratio}
\begin{proof}
\emph{Part~1.} The proof of \cref{thm:p15-conc-benchmarks} gives
\(Ee^{\lambda Z}=e^{\lambda^2/2}\) for standard normal \(Z\). With
\(S_n=\sqrt{n\sigma^2}Z\), \(e^{-hnt}Ee^{hS_n}=\exp(-hnt+h^2n\sigma^2/2)\).
The exponent is a strictly convex quadratic in \(h\), minimized at
\(h=t/\sigma^2>0\), where it equals \(-nt^2/(2\sigma^2)\). Exponential
Markov (\cref{thm:p15-conc-exponential}) with
\(P(S_n\geq nt)=\Phi(-\sqrt nt/\sigma)\) and \(y=\sqrt nt/\sigma\) gives
\(\Phi(-y)\leq e^{-y^2/2}\).

\emph{Part~2.} For \(p\geq1/2\), \(g(p)=1/[2p(1-p)]\), so
\(A_2=e^{-nt^2/(2p(1-p))}\). By part~1 of
\cref{thm:p15-w02b-p15-063-chebyshev-comparison},
\(0<\sigma^2\leq p(1-p)\), so \(nt^2/(2\sigma^2)\geq nt^2/(2p(1-p))\) and
the inequality follows, with equality exactly when \(\sigma^2=p(1-p)\).
The equality case of that part is the stated binomial case.

\emph{Part~3.} Put \(m_i=(a_i+b_i)/2\). Since \(|X_i-m_i|\leq(b_i-a_i)/2\),
\[
 \operatorname{Var}X_i=E(X_i-m_i)^2-(EX_i-m_i)^2\leq\frac{(b_i-a_i)^2}4 .
\]
Summing, \(\sigma^2\leq V_n/(4n)\), so
\(nt^2/(2\sigma^2)\geq 2n^2t^2/V_n\), which is the inequality.
Equality holds exactly when every term attains its bound. For \(a_i=b_i\)
this is automatic. For \(a_i<b_i\) it requires \(EX_i=m_i\) and
\((X_i-m_i)^2=(b_i-a_i)^2/4\) almost surely, the latter because
\((b_i-a_i)^2/4-(X_i-m_i)^2\) is nonnegative with zero expectation. Then
\(X_i\in\{a_i,b_i\}\), and the mean \(m_i\) forces probability \(1/2\) at
each endpoint. Conversely such a law has variance \((b_i-a_i)^2/4\).
The bound \(e^{-2n^2t^2/V_n}\) is \cref{thm:p15-conc-unequal-ranges} at the
threshold \(nt\), the event \(\overline X-E\overline X\geq t\).

\emph{Part~4.} By \cref{thm:p15-conc-benchmarks},
\(\Phi(-y)\leq\phi(y)/y\), hence
\[
 \frac{e^{-y^2/2}}{\Phi(-y)}\geq\frac{y\,e^{-y^2/2}}{\phi(y)}=\sqrt{2\pi}\,y .
\]
For \(y>(2\pi)^{-1/2}\) this exceeds one. For
\(0<y\leq(2\pi)^{-1/2}\), \(e^{-y^2/2}\geq e^{-1/(4\pi)}>1/2\), because
\(1/(4\pi)<1/12<\ln2\); and \(\Phi(-y)<\Phi(0)=1/2\), the normal density
being positive. So the bound is strict at every \(y>0\), and the ratio
tends to infinity. The attainment statements are the endpoint assertions
of \cref{thm:p15-conc-common-range} and
\cref{thm:p15-variance-optimality-boundary}.

\emph{Part~5.} With \(\tau=nt/b\) and \(\lambda=bt/\sigma^2\),
\(\tau\lambda/2=nt^2/(2\sigma^2)\), so \(N=e^{-\tau\lambda/2}\). The middle
two inequalities are \(B_1\leq B_2\leq e^{-\tau h_2(\lambda)}\) from
\cref{thm:p15-variance-tail-chain}. The first is
\(B_1/B_2>e^{-nt^3/(2b\sigma^2)}\) from
\cref{thm:p15-variance-bound-comparison}, combined with \(B_2\geq N\).
For the latter, \(B_2=e^{-\tau h_1(\lambda)}\) and
\(h_1(\lambda)=[(1+\lambda)\ln(1+\lambda)-\lambda]/\lambda\leq\lambda/2\):
the function \(\psi(\lambda)=\lambda+\lambda^2/2-(1+\lambda)\ln(1+\lambda)\)
has \(\psi(0)=0\) and \(\psi'(\lambda)=\lambda-\ln(1+\lambda)\geq0\), so
\((1+\lambda)\ln(1+\lambda)-\lambda\leq\lambda^2/2\). For the last
inequality,
\[
 \frac{e^{-\tau h_2(\lambda)}}N
   =\exp\Bigl\{\tau\Bigl(\frac\lambda2-\frac{\lambda}{2(1+\lambda/3)}\Bigr)\Bigr\}
   =\exp\Bigl\{\frac{\tau\lambda^2}{6(1+\lambda/3)}\Bigr\}
   \leq e^{\tau\lambda^2/6},
\]
and \(\tau\lambda^2/6=nbt^3/(6\sigma^4)\). Both outer exponents tend to
zero when \(nt^3\to0\) with \(b,\sigma^2\) fixed, which for fixed \(n\) is
\(\lambda\to0\).
\end{proof}
\end{annexitem}

\begin{annexitem}{annex:p15-w02b-p15-063-cs-equivalence}{thm:p15-w02b-p15-063-cs-equivalence}{Stopping at the first exclusion}
\begin{proof}
Let \(G=\{\theta\in\mathcal C_n\text{ for every }n\}\), so \(c=P(G)\).
For every random time \(\tau\) with \(P(\tau<\infty)=1\),
\(G\cap\{\tau<\infty\}\subseteq\{\tau<\infty,\theta\in\mathcal C_\tau\}\):
on the left, \(\tau\) equals some finite \(n\) and \(\theta\in\mathcal
C_n\). The coverage event is measurable, being the countable union over
\(n\) of \(\{\tau=n\}\cap\{L_n\leq\theta\leq U_n\}\). Hence the coverage
at \(\tau\) is at least \(P(G)-P(\tau=\infty)=c\).

Conversely let \(\varsigma=\inf\{n\geq1:\theta\notin\mathcal C_n\}\), with
\(\inf\varnothing=\infty\). Since \(\theta\) is a constant under the fixed
\(P\), \(\{\theta\notin\mathcal C_m\}=\{L_m>\theta\}\cup\{U_m<\theta\}
\in\mathcal F_m\), and \(\{\varsigma\leq n\}\) is the union of these events
over \(m\leq n\). Thus \(\varsigma\) is a stopping time, and so is
\(\varsigma_N=\min\{\varsigma,N\}\), which is bounded. If \(\varsigma\leq N\),
then \(\varsigma_N=\varsigma\) and \(\theta\notin\mathcal C_{\varsigma}\);
if \(\varsigma>N\), then \(\varsigma_N=N\) and \(\theta\in\mathcal C_N\).
Hence
\[
 P\{\theta\in\mathcal C_{\varsigma_N}\}=1-P\{\varsigma\leq N\}
   \downarrow1-P\{\varsigma<\infty\}=P(G)=c,
\]
by continuity from below of the increasing events \(\{\varsigma\leq N\}\).
The bounded stopping times form a subclass of the almost surely finite
stopping times, which form a subclass of the almost surely finite random
times. The infimum over each class is at least \(c\) by the first
paragraph and at most \(c\) by the second, so all three equal \(c\). The
equivalences follow by comparing \(c\) with \(1-\delta\).

For the running intersections,
\(\{\theta\in\mathcal C'_n\ \forall n\}=\{\theta\in\mathcal C_m\
\forall m\}\), since each \(\mathcal C'_n\) is the intersection of the
first \(n\) sets. Each \(\mathcal C'_n\) is a closed interval with
endpoints \(\max_{m\leq n}L_m\) and \(\min_{m\leq n}U_m\), which are
\(\mathcal F_n\)-measurable; it is empty when the first exceeds the second.
On \(G\) it contains \(\theta\), so emptiness occurs only off \(G\).
The final sentence of the theorem is part~2 of
\cref{thm:p15-w02b-p15-063-reuse-failure}.
\end{proof}
\end{annexitem}

\begin{annexitem}{annex:p15-w02b-p15-063-union-band-cs}{thm:p15-w02b-p15-063-union-band-cs}{Applying the countable allocation to centered observations}
\begin{proof}
Fix \(P\in\mathcal P\) and write \(\mu=\mu(P)\). Put \(D_i=X_i-\mu\).
Then \(D_i\) is \(\mathcal F_i\)-measurable,
\(E(D_i\mid\mathcal F_{i-1})=0\), and \(a_i-\mu\leq D_i\leq b_i-\mu\) with
deterministic bounds of width \(b_i-a_i\). These are the hypotheses of
\cref{thm:p15-conc-martingale}, with the same \(V_n\), and
\(M_n=\sum_{i\leq n}D_i=n(\overline X_n-\mu)\). By
\cref{thm:p15-conc-time-uniform}, with probability at least \(1-\delta\),
\(|M_n|\leq\rho_n\) for every \(n\geq1\), with radius zero when
\(V_n=0\). Dividing by \(n\), this is the event that
\(\mu\in\mathcal C_n\) for every \(n\). The endpoints are
\(\mathcal F_n\)-measurable, so part~1 holds. Independent identically
distributed observations in \([a,b]\) with natural filtration satisfy
\(E(X_i\mid\mathcal F_{i-1})=EX_1\) by independence of \(X_i\) from
\(\mathcal F_{i-1}\).

For part~2, \(V_n=n\) and
\(\rho_n/n=\sqrt{(n/2)\ln(2n(n+1)/\delta)}/n=r_n\). Since
\(0\leq X_i\leq1\), \(\mu\in[0,1]\); intersecting an interval containing
\(\mu\) with \([0,1]\) keeps \(\mu\), so the clipped intervals have the
same simultaneous coverage event.

For part~3, put \(D_i=X_i-\mu_i\). The same hypotheses hold, now with
deterministic bounds \(a_i-\mu_i,b_i-\mu_i\) of width \(b_i-a_i\), and
\(M_n=n(\overline X_n-\overline\mu_n)\). The same theorem gives
\(|\overline X_n-\overline\mu_n|\leq\rho_n/n\) for every \(n\) with
probability at least \(1-\delta\).
\end{proof}
\end{annexitem}

\begin{annexitem}{annex:p15-w02b-p15-063-reuse-failure}{thm:p15-w02b-p15-063-reuse-failure}{Oscillation of normalized sums and the exact exclusion recursion}
\begin{proof}
As in the preceding proof, \(M_n=n(\overline X_n-\mu)\) has
\(V_n=n\).

\emph{Part~1.} \Cref{thm:p15-conc-martingale} gives
\(P\{|M_n|\geq s\}\leq2e^{-2s^2/n}\) for \(s>0\). At \(s=nf_n\), the
right side is \(2\exp(-2nf_n^2)=2\exp(-\ln(2/\delta))=\delta\). The event
\(\mu\notin\mathcal I^{\rm fix}_n\) is \(|M_n|>nf_n\), contained in the
event just bounded.

\emph{Part~2.} The variables \(Y_i=X_i-\mu\) are independent, identically
distributed and centered, with variance \(\sigma^2\in(0,1/4]\) by part~1
of \cref{thm:p15-w02b-p15-063-chebyshev-comparison}. By
\cref{thm:p15-clt-oscillation},
\(\limsup_n n^{-1/2}\sum_{i\leq n}Y_i=+\infty\) almost surely. Put
\(\kappa_\delta=\sqrt{\ln(2/\delta)/2}\). Then
\(\mu\notin\mathcal I^{\rm fix}_n\) exactly when
\(n^{-1/2}|\sum_{i\leq n}Y_i|>\kappa_\delta\), because
\(nf_n=\sqrt n\,\kappa_\delta\). On the almost sure event where the upper
limit is infinite, the normalized sum exceeds \(\kappa_\delta\) for
infinitely many \(n\). Therefore the simultaneous coverage probability
is zero. By \cref{thm:p15-w02b-p15-063-cs-equivalence} the infimum of the
coverage over bounded stopping times is also zero, attained in the limit
by \(\varsigma_N=\min\{\varsigma,N\}\), where \(\varsigma\) is the first
exclusion time.

\emph{Part~3.} If \(\sigma^2=0\), then \(E(X_1-\mu)^2=0\), so
\(X_1=\mu\) almost surely, and likewise every \(X_i\). Then
\(\overline X_n=\mu\) and each interval contains \(\mu\).

\emph{Part~4.} As in part~1, \(P\{|M_n|>nu_n^{(N)}\}\leq
2\exp(-\ln(2N/\delta))=\delta/N\) for each \(n\leq N\). The finite union
over \(n\leq N\) has probability at most \(\delta\). The radii carry no
information about \(n>N\).

\emph{Part~5.} Write \(T_m=\sum_{i\leq m}X_i\) and
\(G_m=\{p\in\mathcal I^{\rm fix}_j\text{ for }1\leq j\leq m\}\), with
\(G_0=\Omega\). Coverage at time \(j\) is the event \(T_j\in\mathsf A_j\).
We show \(v_m(k)=P\{T_m=k,\ G_m\}\) by induction. At \(m=0\) both sides
are \(\mathbf1_{\{k=0\}}\). For \(m\geq1\),
\[
\begin{split}
 \{T_m=k\}\cap G_m={}&\{k\in\mathsf A_m\}\cap G_{m-1}\\
 &\cap\bigl(\{X_m=1,T_{m-1}=k-1\}\cup\{X_m=0,T_{m-1}=k\}\bigr),
\end{split}
\]
a disjoint union. For each \(j\), the event
\(\{T_{m-1}=j\}\cap G_{m-1}\) is determined by \(X_1,\ldots,X_{m-1}\) and
is therefore independent of \(X_m\). Their probabilities are
\(v_{m-1}(j)\) by the induction hypothesis, and are zero for
\(j\notin\{0,\ldots,m-1\}\). Multiplying by \(P(X_m=1)=p\) and
\(P(X_m=0)=1-p\) gives the recursion. Summing over \(k\) gives
\(P(G_N)\), whose complement is the probability of some exclusion by
time \(N\).
\end{proof}
\end{annexitem}

\begin{annexitem}{annex:p15-w02b-p15-063-allocation}{thm:p15-w02b-p15-063-allocation}{Obtained levels from the fixed-time bound and the first Holm comparison}
\begin{proof}
\emph{Part~1.} For \(u=1\) the inequality is trivial. For \(u=0\) the
event is empty, since \(\pi_n(\mu_0)>0\). For \(0<u<1\), \(\pi_n(\mu_0)\leq u\) is equivalent to
\(2\exp(-2n(\overline X_n-\mu_0)^2)\leq u\), that is to
\(|M_n|\geq ns\) with \(s=\sqrt{\ln(2/u)/(2n)}\) and
\(M_n=n(\overline X_n-\mu_0)\). As in part~1 of the previous proof, its
probability is at most \(2e^{-2ns^2}=u\).

\emph{Part~2.} The map \(d\mapsto2\exp(-2nd^2)\) is strictly decreasing
on \(d\geq0\). Hence \(|\overline X_n-\mu_0|>r_n\) holds exactly when
\(2\exp(-2n(\overline X_n-\mu_0)^2)<2\exp(-2nr_n^2)=\delta/[n(n+1)]=
\delta w_n\). Since \(\delta w_n<1\), the cap at one in \(\pi_n\) does not
affect this comparison. For \(\mu_0\in[0,1]\), clipping to \([0,1]\) does
not change membership. The other two equivalences are the same
computation with \(2\exp(-2nf_n^2)=\delta\) and
\(2\exp(-2n(u^{(N)}_n)^2)=\delta/N\), both less than one.

\emph{Part~3.} The weights satisfy
\(\sum_n\delta w_n=\delta\sum_n(1/n-1/(n+1))=\delta\). By parts~1
and~2 and countable subadditivity, the probability that the true value
is excluded at some time is at most \(\sum_nP\{\pi_n(\mu)<\delta w_n\}
\leq\delta\); the finite-horizon case uses \(N\) terms of \(\delta/N\).
Subadditivity requires nothing about the joint law of the events. The
unadjusted rule excludes \(\mu_0\) whenever some \(\pi_n(\mu_0)<\delta\),
which is the rule \(\mathcal R_{\rm U}\) of
\cref{def:p15-mt-comparison-rules} applied at every time.

\emph{Part~4.} In \cref{def:p15-mt-holm-rule}, the first comparison is
\(R^{(1)}\leq\delta/N\), where \(R^{(1)}=\min_{n\leq N}\pi_n(\mu_0)\). If
it fails, Holm stops with no rejection; if it passes, Holm rejects at
least \(H^{(1)}\). The Bonferroni set
\(\mathcal R_{\rm B}=\{n:\pi_n(\mu_0)\leq\delta/N\}\) is nonempty exactly
when the same minimum is at most \(\delta/N\). For the weighted procedure
of \cref{def:p15-mt-weighted}, the first comparison is
\(S^{(1)}=\min_n\pi_n(\mu_0)/c_n\leq\delta/W_1\) with
\(W_1=\sum_{\ell\leq N}c_\ell\), which holds exactly when some
\(\pi_n(\mu_0)\leq\delta c_n/W_1\). In each case the step-down procedure
rejects some member exactly when the single-step procedure does. Since
every member states \(\mu=\mu_0\), the value \(\mu_0\) is rejected by the
one procedure exactly when it is rejected by the other. These are
pointwise identities and use no distributional assumption.

\emph{Part~5.} Both radii have the form \(\sqrt{\ln(c/\delta)/(2n)}\), with
\(c=2N\) and \(c=2n(n+1)\) respectively, and \(\ln\) is strictly
increasing.
\end{proof}
\end{annexitem}

\begin{annexitem}{annex:p15-w02b-p15-063-sprt-comparison}{thm:p15-w02b-p15-063-sprt-comparison}{Likelihood-ratio martingales and the mean-scale boundary}
\begin{proof}
For \(x\in\{0,1\}\) write \(f_i(x)=p_i^x(1-p_i)^{1-x}\). Both are positive,
so \(L_m=\prod_{j\leq m}f_1(X_j)/f_0(X_j)>0\) and
\(Z_m=\sum_{j\leq m}z_j\) with \(z_j=z_+\) when \(X_j=1\) and \(z_-\) when
\(X_j=0\); \(z_+=\ln(p_1/p_0)>0>z_-=\ln((1-p_1)/(1-p_0))\).

\emph{Part~1.} Under \(P_{p_0}\),
\(E\{f_1(X_j)/f_0(X_j)\}=p_1+(1-p_1)=1\). Since \(X_m\) is independent of
\(\mathcal F_{m-1}=\sigma(X_1,\ldots,X_{m-1})\),
\(E(L_m\mid\mathcal F_{m-1})=L_{m-1}\). So \((L_m)_{m\geq0}\), \(L_0=1\), is a
positive martingale, in particular a submartingale with
\(L_m^+=L_m\) and \(EL_n=1\). By \cref{thm:p15-doob-weak-maximal},
\(A\,P\{\max_{0\leq m\leq n}L_m\geq A\}\leq EL_n=1\). As
\(L_0=1<A\), the event equals \(\{Z_m\geq a\text{ for some }1\leq m\leq
n\}\). These events increase with \(n\) to the event in part~1, and
continuity from below gives the bound \(1/A\).

\emph{Part~2.} With \(T_m=\sum_{j\leq m}X_j=m\overline X_m\),
\(Z_m=T_mz_++(m-T_m)z_-=mz_-+T_m\Delta\). Since \(\Delta>0\),
\(Z_m\geq a\) is equivalent to \(T_m\geq(a-mz_-)/\Delta\), that is to
\(\overline X_m\geq c_*+a/(m\Delta)\). The mean score under
\(\operatorname{Bernoulli}(q)\) is \(e(q)=qz_++(1-q)z_-=z_-+q\Delta\),
strictly increasing in \(q\) with its zero at \(q=c_*\). The ratio
\(f_0/f_1\) takes two distinct values, each with positive probability under
\(P_{p_1}\), so strict concavity of the logarithm gives
\(-e(p_1)=E_{p_1}\ln(f_0/f_1)<\ln E_{p_1}(f_0/f_1)=\ln1=0\). In the same
way \(e(p_0)=E_{p_0}\ln(f_1/f_0)<0\). Hence \(e(p_0)<0<e(p_1)\) and
\(p_0<c_*<p_1\).

\emph{Part~3.} Let \(\varphi(s)=pe^{sz_+}+(1-p)e^{sz_-}\). It is
continuous, \(\varphi(0)=1\), \(\varphi'(0)=e(p)<0\) because \(p<c_*\),
\(\varphi''(s)=pz_+^2e^{sz_+}+(1-p)z_-^2e^{sz_-}>0\), and
\(\varphi(s)\geq pe^{sz_+}\to\infty\) as \(s\to\infty\). So \(\varphi<1\)
just to the right of zero and \(\varphi>1\) for large \(s\); the
intermediate value theorem gives a root \(h>0\). A strictly convex
function takes the value one at most twice, and one of these is \(s=0\);
so \(h\) is unique. Under \(P_p\), \(Y_m=e^{hZ_m}\) satisfies
\(E(Y_m\mid\mathcal F_{m-1})=Y_{m-1}\varphi(h)=Y_{m-1}\), \(Y_0=1\), so it is
a positive martingale of mean one. Since \(Z_m\geq a\) exactly when
\(Y_m\geq A^h\), and \(A^h>1=Y_0\), Doob's inequality and continuity from
below give \(P_p\{Z_m\geq a\text{ for some }m\}\leq A^{-h}<1\).

For the union band, \(\overline X_n\to p\) almost surely by
\cref{thm:p15-iid-series-strong-law}, and \(r_n\to0\) because
\(\ln(2n(n+1)/\delta)/n\to0\). Since \(|p-p_0|>0\), on the almost sure
convergence event \(|\overline X_n-p_0|>r_n\) for all large \(n\), which
by part~2 of \cref{thm:p15-w02b-p15-063-allocation} means
\(p_0\notin\mathcal C_n\).

\emph{Part~4.} If \(p>c_*\), then almost surely
\(\overline X_m\to p\) while \(c_*+a/(m\Delta)\to c_*<p\). So the strict
inequality \(\overline X_m>c_*+a/(m\Delta)\) holds for all large \(m\)
on that event, and by part~2 \(Z_m\geq a\) for such \(m\).

For the remark on the two-boundary test, \cref{thm:p15-se-errors} gives
\(A\alpha'\leq1-\beta'\leq1\).
\end{proof}
\end{annexitem}

\begin{annexitem}{annex:p15-w02b-p15-063-epoch-band}{thm:p15-w02b-p15-063-epoch-band}{One maximal inequality per epoch and a telescoping allocation}
\begin{proof}
For \(k\geq0\) put
\[
 t_k=\sqrt{\frac{V_{N_k}}2\ln\frac{2(k+1)(k+2)}\delta},\qquad
 E_k=\Bigl\{\max_{0\leq j\leq N_k}|M_j|>t_k\Bigr\}.
\] If \(V_{N_k}>0\), then
\(t_k>0\), and \cref{thm:p15-conc-martingale} at the horizon \(N_k\) gives
\[
 P(E_k)\leq P\{\max_{j\leq N_k}|M_j|\geq t_k\}
   \leq2\exp(-2t_k^2/V_{N_k})=\frac{\delta}{(k+1)(k+2)}.
\]
If \(V_{N_k}=0\), then \(t_k=0\) and every \(M_j\), \(j\leq N_k\), is zero
almost surely by the same theorem, so \(P(E_k)=0\). The telescoping sum
\(\sum_{k\geq0}1/[(k+1)(k+2)]=1\) and countable subadditivity give
\(P(\bigcup_kE_k)\leq\delta\).

Off \(\bigcup_kE_k\), every \(n\) lies in its epoch \(k=k(n)\), so
\(n\leq N_k\) and \(|M_n|\leq t_k=\varepsilon_n\) before the zero
replacement. For each \(n\) with \(V_n=0\), \(M_n=0\) almost surely by
\cref{thm:p15-conc-martingale}; the countable union of the null events
\(\{V_n=0,\ M_n\ne0\}\) is null. Off \(\bigcup_kE_k\) and this null set,
\(|M_n|\leq\varepsilon_n\) for every \(n\), including the replaced radii.
This proves the first statement.

In the \([0,1]\) model, apply it to \(D_i=X_i-\mu\) as in the proof of
\cref{thm:p15-w02b-p15-063-union-band-cs}: then \(V_{N_k}=N_k\), \(V_n=n>0\)
and \(\varepsilon_n/n=e_n\). Clipping keeps \(\mu\in[0,1]\), and the
endpoints are \(\mathcal F_n\)-measurable.
\end{proof}
\end{annexitem}

\begin{annexitem}{annex:p15-w02b-p15-063-width-comparison}{thm:p15-w02b-p15-063-width-comparison}{Direct comparison of the closed forms}
\begin{proof}
The two displayed ratios follow by squaring the definitions; the second
exceeds one because \(2n(n+1)\geq4>2\). Write
\(\mathcal N=\ln(2(k+1)(k+2)/\delta)\) and
\(\mathcal D=\ln(2n(n+1)/\delta)\); both are positive since
\(\delta<1\).

\emph{Part~1.} At \(n=1\), \(k=0\), \(N_0=1\) and
\(2(k+1)(k+2)=4=2n(n+1)\), so the ratio is one and
\(e_1=r_1=\sqrt{\ln(4/\delta)/2}\). If \(r_1\geq1\), then
\([X_1-r_1,X_1+r_1]\supseteq[0,1]\) for every \(X_1\in[0,1]\). If \(r_1<1\),
the value \(X_1=0\) gives the clipped interval \([0,r_1]\ne[0,1]\).
And \(r_1\geq1\) exactly when \(\ln(4/\delta)\geq2\), that is
\(\delta\leq4e^{-2}\).

\emph{Part~2.} At \(n=2\), \(k=1\), \(N_1=3\) and \(2(k+1)(k+2)=12=2n(n+1)\),
so the squared ratio is \(3/2\). At \(n=2^k\), \(k\geq1\),
\(N_k/n=2-2^{-k}>1\), while \(\mathcal N/\mathcal D\to1\) as
\(\delta\downarrow0\), since both logarithms are \(\ln(1/\delta)\) plus a
constant. So the ratio exceeds one for small \(\delta\).
At \(n=N_k\), \(k\geq1\), \(N_k/n=1\) and the ratio is
\(\mathcal N/\mathcal D\). By induction \(2^{k+1}\geq k+3\) for \(k\geq1\),
so \(N_k\geq k+2\) and \(N_k+1\geq k+3>k+1\); hence
\(N_k(N_k+1)>(k+1)(k+2)\) and \(\mathcal N<\mathcal D\).

\emph{Part~3.} For \(n\) in epoch \(k\), \(2^k\leq n\leq N_k<2^{k+1}\), so
\(N_k/n<2\) and the squared ratio is less than \(2\mathcal N/\mathcal D\).
This is at most one when \(\mathcal D\geq2\mathcal N\), that is when
\(2n(n+1)/\delta\geq4(k+1)^2(k+2)^2/\delta^2\), equivalently
\(\delta n(n+1)\geq2(k+1)^2(k+2)^2\). Because \(n\geq2^k\), this holds if
\(\delta2^k(2^k+1)\geq2(k+1)^2(k+2)^2\). Call the latter condition
\(\mathsf Q_k\).

If \(\mathsf Q_k\) holds for some \(k\geq2\), it holds for \(k+1\). The
left side is multiplied by
\(2(2^{k+1}+1)/(2^k+1)=2(2-1/(2^k+1))\geq18/5\), and the right side by
\(((k+3)/(k+1))^2\leq(5/3)^2=25/9<18/5\). Hence \(\mathsf Q_k\) holds for
all \(k\geq k_0(\delta)\), and every \(n\geq2^{k_0(\delta)}\) lies in such
an epoch; the ratio is strictly below one there.

For the limit, \(k\leq\log_2n\) gives
\(\mathcal N\leq\ln(2(\log_2n+1)(\log_2n+2)/\delta)\), whose quotient by
\(\ln n\) tends to zero, while \(\mathcal D\geq\ln(2n^2)\geq2\ln n\). So
\(2\mathcal N/\mathcal D\to0\).

The values of \(k_0\) follow from the table of both sides of
\(\mathsf Q_k\), the left side at \(\delta=0.05\):
\[
\begin{array}{r|cccccccc}
 k&2&3&4&5&6&7&8&9\\ \hline
 2^k(2^k+1)&20&72&272&1056&4160&16512&65792&262656\\
 0.05\cdot2^k(2^k+1)&1&3.6&13.6&52.8&208&825.6&3289.6&13132.8\\
 2(k+1)^2(k+2)^2&288&800&1800&3528&6272&10368&16200&24200
\end{array}
\]
At \(\delta=0.05\), \(\mathsf Q_k\) fails for \(2\leq k\leq9\), while
\(\mathsf Q_{10}\) holds: \(0.05\cdot1024\cdot1025=52480\geq34848=
2\cdot11^2\cdot12^2\). So \(k_0(0.05)=10\). For every \(\delta\leq1\),
\(\mathsf Q_k\) fails for \(2\leq k\leq6\), and \(\mathsf Q_7\) reads
\(16512\,\delta\geq10368\).

\emph{Part~4.} Each radius is \(\sqrt{C\ln(c/\delta)}\) with constants
\(C>0\) and \(c\in\{2n(n+1),2(k+1)(k+2),2\}\), each greater than one. As
\(\delta\uparrow1\), \(\ln(c/\delta)\) decreases to \(\ln c>0\).

\emph{Part~5.} \(\rho_n=0\) exactly when \(V_n=0\) by definition. Before
replacement, \(\varepsilon_n\) is positive exactly when \(V_{N_k}>0\).
That the replacement changes the band only on a null event was shown in
the proof of \cref{thm:p15-w02b-p15-063-epoch-band}. For widths
\(2^i\), \(V_n=\sum_{i\leq n}4^i=(4^{n+1}-4)/3\). At \(n=2^k\),
\(k\geq1\), \(N_k=2n-1\), so
\(V_{N_k}/V_n\geq4^{N_k}/(4^{n+1}/3)=3\cdot4^{2^k-2}\). Also
\(\mathcal N\geq\ln4\), and \(2n(n+1)\leq4n^2=4^{k+1}\) gives
\(\mathcal D\leq\ln(4^{k+1}/\delta)\). Therefore
\[
 \frac{\varepsilon_n^2}{\rho_n^2}
   =\frac{V_{N_k}}{V_n}\cdot\frac{\mathcal N}{\mathcal D}
   \geq\frac{3\cdot4^{2^k-2}\ln4}{(k+1)\ln4+\ln(1/\delta)}
   \longrightarrow\infty\qquad(k\to\infty).
\]
\end{proof}
\end{annexitem}
